\documentclass[10pt,a4paper]{amsart}
\usepackage[margin=19mm]{geometry}
\usepackage{amsmath,amssymb,mathtools}
\usepackage[scr=rsfs,cal=euler]{mathalfa}
\usepackage{xfrac}
\usepackage[unicode,hidelinks,bookmarks=false]{hyperref}
\newcommand{\ie}{i.e.\ }
\newcommand{\cf}{cf.\ }
\newcommand{\resp}{respectively\ }
\newcommand{\Subsection}{\subsection}
\providecommand{\nobreakdash}{\nobreak\hbox{-}}
\setlength{\emergencystretch}{2em}
\multlinegap0pt
\usepackage{amssymb}
\newtheorem{theorem}{Theorem}
\title{A finite-sample Borel--Cantelli inequality under $m$-dependence}
\author{Chatchawan Panraksa}
\date{Source: arXiv:2604.07750v1 (CC BY 4.0)}
\begin{document}
\maketitle

\noindent\emph{Source and licence.} A finite-sample Borel--Cantelli inequality under $m$-dependence. Version arXiv:2604.07750v1 (CC BY 4.0), distributed by arXiv under the Creative Commons Attribution 4.0 International licence, http://creativecommons.org/licenses/by/4.0/, as recorded in the arXiv licence metadata for this exact version and shown on the abstract page https://arxiv.org/abs/2604.07750v1. Full licence notice: \texttt{licenses/arxiv-2604.07750-CC-BY-4.0.txt}.\par\smallskip

\noindent\textit{Editorial scope.} Counted unit: the article's theorem thm:cov (source0.tex lines 321--332). The article's complete original proof of it is delivered with it (source0.tex lines 334--449, one \texttt{proof} environment of 116 lines). No mathematics is written, repaired or re-derived: the statement and the proof are the article's own lines, in the article's own notation, and no retained line is substituted or re-flowed.  Retained uncounted as the source-local context the proof needs: lines 177--189.  Nothing was omitted from any retained span, so the package carries no omission record.  No retained line contains a citation, so the wrapper carries no bibliography and no bibliography entry.  No retained line contains a figure environment, an image-inclusion command or drawing code, so no image is delivered and the package declares no asset dependency.  Editorial formatting only, in addition: an AMS \texttt{amsart} preamble that restates the article's own declarations and macros used by the retained lines, namely the environments \texttt{theorem} on separate counters, together with the standard packages those lines need.  The retained environments print in this document as Theorem 1 in order of appearance, whereas the article prints the same items with the numbering of its own full document.  The complete native source of the article as distributed in the arXiv e-print for this exact version is bundled unchanged as \texttt{sources/198113/source0.tex}.
\begin{theorem}[Finite-sample $m$-dependence inequality]
\label{thm:finite}
Let $m\ge 1$ be an integer and let $(A_k)_{1\le k\le N}$ be a finite
sequence of $m$-dependent events.  Write
\[
  S_N=\sum_{k=1}^{N} \mathbb{P}(A_k).
\]
Then
\[
  \mathbb{P}\Bigl(\bigcup_{k=1}^N A_k\Bigr)
  \ge 1 - \exp\Bigl(- \frac{1}{m+1}\,S_N\Bigr).
\]
\end{theorem}

\begin{theorem}[Local-intersection inequality]\label{thm:cov}
Under the assumptions of Theorem~\ref{thm:finite}, we have
\begin{equation}\label{eq:cov-refined}
  \mathbb{P}\Bigl(\bigcup_{k=1}^N A_k\Bigr)
  \;\ge\;
  1-\exp\Biggl(
   -\frac{1}{2}\sum_{k=1}^N \mathbb{P}(A_k)
   \;+\;\frac{1}{2}\!\!\sum_{\substack{1\le i<j\le N\\|i-j|\le m-1}}
   \mathbb{P}(A_i\cap A_j)
  \Biggr).
\end{equation}
\end{theorem}

\begin{proof}
\textit{Step 1 (Shifted block partitions and $1$-dependence).}\par
For each $r\in\{0,1,\dots,m-1\}$, define the $r$-shifted partition of
$\{1,\dots,N\}$ into disjoint length-$m$ blocks by
\[
\begin{aligned}
  I^{(r)}_{j}
  ={}&\{\,r+(j-1)m+1,\ r+(j-1)m+2,\ \dots,\ r+jm\,\}\\
      &{}\cap\{1,\dots,N\}, \qquad j=0,1,2,\dots,
\end{aligned}
\]
and the corresponding block events
\[
  B^{(r)}_{j}=\bigcup_{k\in I^{(r)}_{j}} A_k.
\]
For $j=0$ the block may be shorter; the sets $(I^{(r)}_{j})_{j\ge0}$ are
disjoint and their union is $\{1,\dots,N\}$.  Moreover, the block sequence
$\bigl(B^{(r)}_{j}\bigr)_{j\ge0}$ is $1$-dependent: if $|j-j'|\ge2$ and
$k\in I^{(r)}_{j}$, $k'\in I^{(r)}_{j'}$, then $|k-k'|\ge m+1>m$, so
$B^{(r)}_{j}$ and $B^{(r)}_{j'}$ are independent by $m$-dependence.

\medskip
\textit{Step 2 (Parity splitting and product-to-exponential).}
Fix $r$.  Let $\mathcal{J}_{\mathrm{odd}}^{(r)}$ and
$\mathcal{J}_{\mathrm{even}}^{(r)}$ be the odd/even block indices.  Since the
blocks within each parity are mutually independent and
\[
  \bigcup_{k=1}^N A_k = \bigcup_{j\ge0} B^{(r)}_{j},
\]
we have
\begin{align*}
\mathbb{P}\Bigl(\bigcap_{k=1}^N A_k^{\mathrm c}\Bigr)
  &= \mathbb{P}\Bigl(\bigcap_{j\ge0}(B^{(r)}_{j})^{\mathrm c}\Bigr)\\
  &\le \min\!\Bigl\{
  \prod_{j\in\mathcal{J}_{\mathrm{odd}}^{(r)}}\bigl(1-\mathbb{P}(B^{(r)}_{j})\bigr),\,
  \prod_{j\in\mathcal{J}_{\mathrm{even}}^{(r)}}\bigl(1-\mathbb{P}(B^{(r)}_{j})\bigr)
  \Bigr\}.
\end{align*}
Using $1-x\le e^{-x}$ and $\min\{e^{-x},e^{-y}\}\le e^{-(x+y)/2}$, we obtain
\begin{equation}\label{eq:parity-exp}
  \mathbb{P}\Bigl(\bigcup_{k=1}^N A_k\Bigr)
  \;\ge\; 1-\exp\Bigl(-\frac{1}{2}\sum_{j\ge0}\mathbb{P}(B^{(r)}_{j})\Bigr)
  \qquad\text{for each }r\in\{0,1,\dots,m-1\}.
\end{equation}

\medskip
\textit{Step 3 (Second-order Bonferroni inside blocks).}
Bonferroni’s second-order inequality yields
\[
  \mathbb{P}\bigl(B^{(r)}_{j}\bigr)
  \ge
  \sum_{i\in I^{(r)}_{j}} \mathbb{P}(A_i)
  -
  \sum_{\substack{i<\ell\\ i,\ell\in I^{(r)}_{j}}} \mathbb{P}(A_i\cap A_\ell).
\]
Summing over $j\ge0$ and averaging over $r=0,\dots,m-1$ gives
\begin{equation}\label{eq:block-avg}
  \frac{1}{m}\sum_{r=0}^{m-1}\ \sum_{j\ge0}\ \mathbb{P}\bigl(B^{(r)}_{j}\bigr)
  \ge
  \sum_{k=1}^N \mathbb{P}(A_k)
  -
  \frac{1}{m}\sum_{r=0}^{m-1}\ \sum_{j\ge0}
  \sum_{\substack{i<\ell\\ i,\ell\in I^{(r)}_{j}}} \mathbb{P}(A_i\cap A_\ell).
\end{equation}

\medskip
\textit{Step 4 (Pair-counting across shifted partitions).}
Fix $1\le i<\ell\le N$ with gap $d=\ell-i$.  The pair $(i,\ell)$ can lie in
a common length-$m$ block only when $d\le m-1$, and in that case this occurs
for exactly $m-d$ values of $r\in\{0,\dots,m-1\}$.  Consequently,
\[
  \frac{1}{m}\sum_{r=0}^{m-1}\ \sum_{j\ge0}
  \sum_{\substack{i<\ell\\ i,\ell\in I^{(r)}_{j}}} \mathbb{P}(A_i\cap A_\ell)
  \le
  \sum_{\substack{1\le i<\ell\le N\\ |i-\ell|\le m-1}} \mathbb{P}(A_i\cap A_\ell).
\]
Plugging this into \eqref{eq:block-avg} yields
\begin{equation}\label{eq:LB-sumB}
  \frac{1}{m}\sum_{r=0}^{m-1}\ \sum_{j\ge0}\ \mathbb{P}\bigl(B^{(r)}_{j}\bigr)
  \ge
  \sum_{k=1}^N \mathbb{P}(A_k)
  -
  \sum_{\substack{1\le i<\ell\le N\\ |i-\ell|\le m-1}} \mathbb{P}(A_i\cap A_\ell).
\end{equation}

\medskip
\textit{Step 5 (From block sums to the exponential bound).}
Let
\[
  X_r:=\sum_{j\ge0}\mathbb{P}(B^{(r)}_{j}).
\]
From \eqref{eq:parity-exp} we have
\[
  \mathbb{P}\Bigl(\bigcup_{k=1}^N A_k\Bigr)
  \ge 1-\exp\Bigl(-\frac{1}{2}X_r\Bigr)
  \qquad\text{for each }r.
\]
Hence
\[
  \mathbb{P}\Bigl(\bigcup_{k=1}^N A_k\Bigr)
  \ge 1-\exp\Bigl(-\frac{1}{2}\max_{0\le r\le m-1} X_r\Bigr)
  \ge 1-\exp\Bigl(-\frac{1}{2m}\sum_{r=0}^{m-1} X_r\Bigr).
\]
Using \eqref{eq:LB-sumB} to lower bound
\(\frac{1}{m}\sum_{r=0}^{m-1} X_r\), we obtain
\[
  \mathbb{P}\Bigl(\bigcup_{k=1}^N A_k\Bigr)
  \ge
  1-\exp\Biggl(
   -\frac{1}{2}\sum_{k=1}^N \mathbb{P}(A_k)
   +\frac{1}{2}\!\!\sum_{\substack{1\le i<\ell\le N\\|i-\ell|\le m-1}}
   \mathbb{P}(A_i\cap A_\ell)
  \Biggr),
\]
which is exactly \eqref{eq:cov-refined}.
\end{proof}

\end{document}
